\section{Introduction}

The big red button is already overdescribed before anyone enters the laboratory. It appears as emergency stop, launch switch, comic foreshadowing device, and meme of irresistible action. What unifies these scenes is not merely redness or scale, but a peculiar invitation: the button is legible as an event waiting to happen. Histories of pushing show that buttons have long condensed pleasure, panic, command, and delegated force into a single tactile interface \citep{Plotnick2018,PlotnickInterface2012,PlotnickWarfare2012}. The red button, in particular, acquires cultural authority by staging action as both immediate and consequential.

This matters for HCI because buttons are not neutral geometric primitives. They are design artifacts that compress a wide field of deliberation into a discrete, countable act. Affordance accounts describe how artifacts invite action \citep{Gaver1991,Norman1999}; experience-centered HCI further reminds us that interaction inherits mood, memory, expectation, and narrative texture from the situation in which it appears \citep{McCarthyWright2004}. A big red button therefore does more than offer a control surface. It recruits a history of command, risk, and anticipation into the local moment of possible contact.

The present study approaches pressing as a social and psychological problem as well as a technical one. Open-ended situations can make people curious, restless, or eager to transform inaction into an event \citep{Wilson2014,KiddHayden2015}. The meaning of action also changes under observation, instrumentation, and feedback \citep{HallHenningsen2008,HaggardChambon2012}. In this project, participants are not instructed to complete a task; they are told that pressing and non-pressing are both valid outcomes. The question is therefore not whether they can press, but how a deliberately overdetermined button becomes pressable in the first place.

Pressability is also shaped by bodily placement and perceptual salience. In immersive displays, object position affects comfort, reach, and interaction efficiency \citep{Lin2021,Mirbagheri2024,Alghofaili2019,McNamara2019}. The button used here was designed to remain visible, reachable, and chromatically prominent, drawing on ergonomic reach estimates and salience research \citep{Bonin2022,Li2008,CastellottiEtAl2022}. At the same time, the paper does not assume that red has a universal meaning. Color-emotion associations vary across cultures, even when red often carries strong learned links to warning, arousal, or urgency in Western interface traditions \citep{ElliotMaier2014,GaoEtAl2007}. Any account of pressability therefore has to balance the specificity of this object with the limits of its symbolism.

The repository's study materials point toward a gap in the existing literature. Prior work can explain aspects of button salience, affordance, presence, biofeedback, peripersonal space, agency, boredom, and curiosity, but it rarely studies voluntary pressing as a single multilevel phenomenon. The Big Red Button Institute's VR paradigm is designed precisely around that gap: a neutral room, one large red button, explicit permission to press or abstain, a visible cumulative counter, and physiological coupling through live or sham cardiac feedback plus respiration-linked environmental light. Rather than treating symbolic history as decoration and embodiment as mechanism, the study puts them in the same frame.

This paper therefore makes three grounded contributions. First, it operationalizes \emph{pressability} as a multidisciplinary construct that links interface history, perceptual invitation, social observation, and embodied coupling. Second, it documents a compile-ready version of the repository's first VR study protocol, including its placement rationale and biofeedback logic. Third, it states a truthful empirical status: the protocol is well enough specified to write, critique, and refine, but the current repository snapshot does not yet justify reporting participant outcomes. The working hypotheses are correspondingly modest. Live physiological coupling may increase pressing, presence, or felt proximity to the button; novelty and order effects may dominate those condition differences; and null results would still be theoretically meaningful if symbolic and ergonomic saturation already make the button difficult to ignore.
